%% ========================================================================
%% Chapter X: Bash Shell dan Shell Basic
%% ========================================================================

\chapter{Bash Shell dan Shell Basic}
\label{ch:bash-shell-basic}

Bab ini membahas Bash sebagai shell default yang umum digunakan pada Linux,
cara kerja file konfigurasi Bash, variabel lingkungan dan \texttt{PATH}, pembuatan
alias dan fungsi shell, fitur completion dan history, wildcard dan ekspansi nama file,
serta quoting dan escaping. Fokus pembahasan diarahkan pada kebutuhan praktikum
dan administrasi sistem sehari-hari: bekerja lebih cepat di terminal, mengurangi
kesalahan saat mengetik perintah, dan membangun kebiasaan kerja yang aman.

\begin{examplebox}[Bash sebagai ``panel kendali'' administrator]
Bayangkan Anda sedang mengelola sebuah server tanpa antarmuka grafis.
Tidak ada tombol klik, tidak ada menu visual, dan tidak ada jendela aplikasi.
Semua interaksi dilakukan melalui terminal. Dalam kondisi seperti ini,
\textbf{shell} berperan sebagai panel kendali utama antara pengguna dan sistem operasi.

Bash menerima perintah yang Anda ketik, menjalankan program, mengatur
variabel lingkungan, menangani redirection, membantu pencarian perintah
melalui history, bahkan memungkinkan Anda membuat shortcut kerja sendiri
melalui alias dan fungsi. Karena itu, memahami Bash bukan hanya soal
``bisa mengetik command'', tetapi juga soal membangun workflow kerja yang
lebih cepat, konsisten, dan aman.
\end{examplebox}

%% ------------------------------------------------------------------------
\section{Pengenalan Bash sebagai Shell Default di Linux}
\label{sec:pengenalan-bash}

Bash (\textit{Bourne Again Shell}) adalah shell yang sangat umum dijumpai
pada sistem Linux modern. Shell ini menjadi antarmuka yang menerjemahkan
perintah dari pengguna menjadi aksi yang dapat dijalankan sistem operasi.
Bash juga menyediakan banyak fasilitas praktis seperti history, completion,
variabel lingkungan, alias, fungsi, dan scripting sederhana.

Dalam konteks administrasi sistem, Bash penting karena hampir seluruh
aktivitas operasional pada server Linux dilakukan dari terminal:
membaca log, mengelola service, memeriksa disk, membuat backup, hingga
menjalankan otomasi sederhana. Jika terminal adalah ``ruang kerja'' Anda,
maka Bash adalah alat kerja utamanya.

\subsection*{Fungsi Perintah yang Digunakan}

\begin{itemize}
  \item \texttt{echo \$SHELL}: melihat shell login pengguna.
  \item \texttt{echo \$0}: melihat nama shell aktif saat ini.
  \item \texttt{ps -p \$\$ -o pid,ppid,args=}: melihat proses shell yang sedang aktif.
  \item \texttt{bash --version}: melihat versi Bash.
  \item \texttt{help}: bantuan untuk \textit{builtin command}.
  \item \texttt{man bash}: dokumentasi lengkap Bash.
  \item \texttt{mkdir}, \texttt{touch}: menyiapkan workspace praktikum.
\end{itemize}

\begin{tipbox}
Pada Linux, istilah \textit{default shell} bisa merujuk ke shell login pengguna,
sedangkan shell yang sedang aktif pada terminal saat ini bisa saja berbeda.
Karena itu, biasakan membedakan \texttt{\$SHELL} dan \texttt{\$0}.
\end{tipbox}

\subsection*{Praktikum 6.1: Mengenali Bash dan Menyiapkan Workspace }

Tujuan: mengenali shell aktif dan menyiapkan area kerja yang aman untuk seluruh praktikum bab ini.

\begin{enumerate}
  \item Lihat shell login dan shell aktif saat ini:
\begin{lstlisting}[language=bash, caption={Melihat shell login dan shell aktif}]
echo "Shell login : $SHELL"
echo "Shell aktif : $0"
bash --version | head -n 1
\end{lstlisting}

  \item Lihat proses shell yang sedang berjalan:
\begin{lstlisting}[language=bash, caption={Melihat proses shell aktif}]
echo $$
ps -p $$ -o pid,ppid,args=
\end{lstlisting}

  \item Buat workspace praktikum:
\begin{lstlisting}[language=bash, caption={Membuat workspace Bash}]
mkdir -p ~/praktikum-os/week07-bash/{bin,backup,logs,sampel,ruang-nama}
cd ~/praktikum-os/week04-bash
pwd
\end{lstlisting}

  \item Buat beberapa file contoh yang akan dipakai pada praktikum berikutnya:
\begin{lstlisting}[language=bash, caption={Membuat file contoh untuk praktikum lanjutan}]
touch sample-app.conf
touch logs/app-01.log logs/app-02.log logs/app-03.log
touch sampel/catatan-a.txt sampel/catatan-b.txt
touch sampel/backup-01.tar sampel/backup-02.tar
touch sampel/laporan-harian.log sampel/laporan-mingguan.log sampel/laporan-bulanan.log
touch "ruang-nama/laporan server april.txt"
touch "ruang-nama/backup [mingguan] server.conf"
ls -R
\end{lstlisting}
\end{enumerate}

\subsection*{Praktikum 6.2: Membuat Ringkasan Sesi Terminal }

Tujuan: membiasakan administrator memverifikasi konteks kerja sebelum melakukan maintenance.

\begin{enumerate}
  \item Masuk ke workspace praktikum:
\begin{lstlisting}[language=bash, caption={Masuk ke workspace praktikum}]
cd ~/praktikum-os/week04-bash
\end{lstlisting}

  \item Simpan informasi sesi terminal ke file laporan:
\begin{lstlisting}[language=bash, caption={Membuat ringkasan sesi terminal}]
{
  echo "=== RINGKASAN SESI BASH ==="
  date
  echo "User       : $(whoami)"
  echo "Hostname   : $(hostname)"
  echo "Shell login: $SHELL"
  echo "Shell aktif: $0"
  echo "PID shell  : $$"
  echo "Direktori  : $(pwd)"
} | tee session-info.txt
\end{lstlisting}

  \item Verifikasi isi file laporan:
\begin{lstlisting}[language=bash, caption={Membaca ringkasan sesi}]
cat session-info.txt
\end{lstlisting}
\end{enumerate}

\begin{notebox}
Kebiasaan kecil seperti mencatat user, host, shell, dan direktori aktif
sering membantu mencegah kesalahan fatal, misalnya salah bekerja di server
atau direktori yang bukan target maintenance.
\end{notebox}

%% ------------------------------------------------------------------------
\section{Konfigurasi Bash (\texttt{.bashrc} dan \texttt{.bash\_profile})}
\label{sec:konfigurasi-bash}

Bash memiliki beberapa file konfigurasi yang dibaca dalam konteks berbeda.
Secara sederhana, \texttt{.bashrc} biasa dipakai untuk pengaturan shell interaktif:
alias, prompt, fungsi, dan variabel yang ingin tersedia saat membuka terminal.
Sementara itu, \texttt{.bash\_profile} umumnya dipakai untuk shell login.

Pada banyak sistem Ubuntu, file login yang aktif sering kali adalah \texttt{.profile}.
Namun, jika \texttt{.bash\_profile} ada, maka file ini biasanya diprioritaskan oleh Bash.
Karena itu, praktik aman yang umum dilakukan adalah membuat \texttt{.bash\_profile}
agar memanggil \texttt{.bashrc}, sehingga konfigurasi interaktif tetap konsisten.

\subsection*{Fungsi Perintah yang Digunakan}

\begin{itemize}
  \item \texttt{ls -la \textasciitilde}: melihat file konfigurasi pada home directory.
  \item \texttt{cp}: membuat backup file konfigurasi.
  \item \texttt{cat <<'EOF' >> file}: menambahkan blok konfigurasi secara aman.
  \item \texttt{source \textasciitilde/.bashrc}: menerapkan konfigurasi tanpa logout.
  \item \texttt{bash -l}: menjalankan shell login baru untuk pengujian.
\end{itemize}

\begin{notebox}
Jangan mengganti seluruh isi \texttt{.bashrc} atau \texttt{.bash\_profile}
di server produksi tanpa backup. Untuk praktikum, biasakan menambahkan blok
konfigurasi sendiri dan menandainya dengan komentar yang jelas.
\end{notebox}

\subsection*{Praktikum 6.3: Menambahkan Konfigurasi Aman pada \texttt{.bashrc} }

Tujuan: memahami peran \texttt{.bashrc} dan menerapkan perubahan secara aman.

\begin{enumerate}
  \item Lihat file konfigurasi Bash pada home directory:
\begin{lstlisting}[language=bash, caption={Melihat file konfigurasi Bash}]
ls -la ~ | grep -E 'bashrc|bash_profile|profile'
\end{lstlisting}

  \item Buat backup \texttt{.bashrc}:
\begin{lstlisting}[language=bash, caption={Backup .bashrc}]
cp ~/.bashrc ~/.bashrc.bak-praktikum
\end{lstlisting}

  \item Tambahkan blok konfigurasi praktikum:
\begin{lstlisting}[language=bash, caption={Menambahkan konfigurasi ke .bashrc}]
cat <<'EOF' >> ~/.bashrc

#: Praktikum Bash Shell ---
export PRAKTIKUM_BASH_DIR="$HOME/praktikum-os/week04-bash"
export EDITOR=nano
#: End Praktikum Bash Shell ---

EOF
\end{lstlisting}

  \item Terapkan konfigurasi tanpa logout:
\begin{lstlisting}[language=bash, caption={Memuat ulang .bashrc}]
source ~/.bashrc
echo "$PRAKTIKUM_BASH_DIR"
echo "$EDITOR"
\end{lstlisting}
\end{enumerate}

\subsection*{Praktikum 6.4: Menyiapkan \texttt{.bash\_profile} untuk Shell Login }

Tujuan: memahami hubungan \texttt{.bash\_profile} dan \texttt{.bashrc} pada skenario login shell.

\begin{enumerate}
  \item Backup \texttt{.bash\_profile} jika sudah ada:
\begin{lstlisting}[language=bash, caption={Backup .bash\_profile jika tersedia}]
[ -f ~/.bash_profile ] && cp ~/.bash_profile ~/.bash_profile.bak-praktikum
\end{lstlisting}

  \item Tambahkan konfigurasi login shell:
\begin{lstlisting}[language=bash, caption={Menambahkan konfigurasi ke .bash\_profile}]
cat <<'EOF' >> ~/.bash_profile

#: Praktikum Bash Login Shell ---
if [ -f ~/.bashrc ]; then
  . ~/.bashrc
fi

echo "Login Bash pada $(date '+%F %T')" >> "$HOME/praktikum-os/week07-bash/login-audit.log"
#: End Praktikum Bash Login Shell ---

EOF
\end{lstlisting}

  \item Uji dengan membuka login shell baru:
\begin{lstlisting}[language=bash, caption={Menguji .bash\_profile dengan login shell}]
bash -l
tail -n 3 ~/praktikum-os/week07-bash/login-audit.log
exit
\end{lstlisting}
\end{enumerate}

%% ------------------------------------------------------------------------
\section{Variabel Lingkungan dan \texttt{PATH}}
\label{sec:variabel-lingkungan-path}

Variabel shell adalah pasangan nama dan nilai yang disimpan di dalam shell.
Sebagian variabel hanya berlaku lokal pada shell saat ini, sedangkan variabel
lingkungan (\textit{environment variable}) diekspor sehingga ikut diwariskan
ke proses anak. Contoh yang sangat penting bagi administrator adalah
\texttt{PATH}, yaitu daftar direktori yang dicari Bash saat Anda mengetik nama
perintah tanpa path lengkap.

Memahami \texttt{PATH} sangat penting. Jika sebuah script berada di direktori
yang belum masuk \texttt{PATH}, maka script itu tidak bisa dipanggil langsung.
Sebaliknya, konfigurasi \texttt{PATH} yang buruk juga berbahaya, misalnya
menaruh direktori yang tidak terpercaya di depan path sistem.

\subsection*{Fungsi Perintah yang Digunakan}

\begin{itemize}
  \item \texttt{echo \$NAMA\_VAR}: melihat isi variabel.
  \item \texttt{VAR=nilai}: membuat variabel shell lokal.
  \item \texttt{export VAR=nilai}: membuat variabel lingkungan.
  \item \texttt{printenv}: melihat environment variable.
  \item \texttt{which} dan \texttt{type}: melihat lokasi/jenis perintah.
  \item \texttt{chmod +x}: memberi izin eksekusi pada script.
\end{itemize}

\begin{tipbox}
Ingat perbedaannya:
\begin{itemize}
  \item \texttt{VAR=nilai} hanya berlaku di shell saat ini.
  \item \texttt{export VAR=nilai} mewariskan variabel itu ke proses anak.
\end{itemize}
\end{tipbox}

\subsection*{Praktikum 6.5: Membedakan Variabel Shell dan Environment Variable }

Tujuan: memahami perbedaan variabel lokal dan variabel lingkungan.

\begin{enumerate}
  \item Buat variabel lokal:
\begin{lstlisting}[language=bash, caption={Membuat variabel shell lokal}]
KELAS_OS="Sistem Operasi A"
echo "$KELAS_OS"
\end{lstlisting}

  \item Buka subshell dan cek apakah variabel masih ada:
\begin{lstlisting}[language=bash, caption={Menguji variabel lokal pada subshell}]
bash
echo "$KELAS_OS"
exit
\end{lstlisting}

  \item Sekarang ubah menjadi environment variable:
\begin{lstlisting}[language=bash, caption={Membuat environment variable}]
export KELAS_OS="Sistem Operasi A"
bash
echo "$KELAS_OS"
exit
\end{lstlisting}

  \item Lihat isi \texttt{PATH} dan lokasi beberapa perintah:
\begin{lstlisting}[language=bash, caption={Melihat PATH dan lokasi perintah}]
echo "$PATH"
which bash
type ls
\end{lstlisting}
\end{enumerate}

\subsection*{Praktikum 6.6: Menambahkan Direktori Script Pribadi ke \texttt{PATH} }

Tujuan: membuat perintah administrasi sederhana yang bisa dipanggil dari mana saja.

\begin{enumerate}
  \item Pastikan direktori \texttt{bin} praktikum tersedia:
\begin{lstlisting}[language=bash, caption={Menyiapkan direktori bin}]
mkdir -p ~/praktikum-os/week07-bash/bin
\end{lstlisting}

  \item Tambahkan direktori tersebut ke \texttt{PATH} melalui \texttt{.bashrc}:
\begin{lstlisting}[language=bash, caption={Menambahkan direktori bin ke PATH}]
cat <<'EOF' >> ~/.bashrc

#: Praktikum PATH ---
export PATH="$HOME/praktikum-os/week07-bash/bin:$PATH"
#: End Praktikum PATH ---

EOF

source ~/.bashrc
echo "$PATH"
\end{lstlisting}

  \item Buat script ringkasan sistem:
\begin{lstlisting}[language=bash, caption={Membuat script ringkas-sistem}]
cat <<'EOF' > ~/praktikum-os/week07-bash/bin/ringkas-sistem
#!/usr/bin/env bash
echo "Hostname : $(hostname)"
echo "User     : $(whoami)"
echo "Uptime   : $(uptime -p)"
echo "Disk /   :"
df -h /
EOF

chmod +x ~/praktikum-os/week07-bash/bin/ringkas-sistem
\end{lstlisting}

  \item Jalankan script dari direktori yang berbeda:
\begin{lstlisting}[language=bash, caption={Menjalankan script dari sembarang lokasi}]
cd ~
ringkas-sistem
\end{lstlisting}
\end{enumerate}

\begin{notebox}
Menaruh direktori personal di awal \texttt{PATH} memudahkan pemanggilan script,
tetapi tetap harus hati-hati agar nama script Anda tidak bentrok dengan perintah
sistem seperti \texttt{ls}, \texttt{cp}, atau \texttt{find}.
\end{notebox}

%% ------------------------------------------------------------------------
\section{Membuat Alias dan Fungsi Shell}
\label{sec:alias-fungsi-shell}

Saat bekerja di terminal, sering ada perintah yang berulang-ulang dipakai.
Alias membantu mempersingkat perintah yang sederhana, sedangkan fungsi shell
digunakan jika logika yang dibutuhkan lebih kompleks, misalnya menerima argumen,
menjalankan beberapa langkah, atau memuat validasi.

Dalam pekerjaan administrator, alias cocok untuk shortcut yang aman dan sering dipakai,
misalnya \texttt{ll}, \texttt{hist10}, atau \texttt{grep --color=auto}.
Untuk tugas yang lebih ``hidup'', seperti backup file konfigurasi dengan timestamp,
fungsi shell lebih tepat digunakan daripada alias.

\subsection*{Fungsi Perintah yang Digunakan}

\begin{itemize}
  \item \texttt{alias}: membuat dan melihat alias.
  \item \texttt{unalias}: menghapus alias.
  \item \texttt{function\_name() \{ ... \}}: membuat fungsi shell.
  \item \texttt{type}: membedakan apakah sesuatu adalah alias, fungsi, builtin, atau file.
  \item \texttt{source ~/.bashrc}: memuat ulang alias dan fungsi.
\end{itemize}

\subsection*{Praktikum 6.7: Membuat Alias Produktivitas Dasar }

Tujuan: membuat shortcut perintah yang sering dipakai.

\begin{enumerate}
  \item Tambahkan alias ke \texttt{.bashrc}:
\begin{lstlisting}[language=bash, caption={Menambahkan alias ke .bashrc}]
cat <<'EOF' >> ~/.bashrc

#: Praktikum Alias ---
alias ll='ls -lah --color=auto'
alias hist10='history | tail -n 10'
alias cdbashlab='cd $HOME/praktikum-os/week04-bash'
#: End Praktikum Alias ---

EOF

source ~/.bashrc
\end{lstlisting}

  \item Uji alias:
\begin{lstlisting}[language=bash, caption={Menguji alias}]
ll
hist10
cdbashlab
pwd
type ll
\end{lstlisting}
\end{enumerate}

\subsection*{Praktikum 6.8: Membuat Fungsi Backup Konfigurasi }

Tujuan: membuat fungsi shell yang dapat dipakai berulang untuk backup file konfigurasi.

\begin{enumerate}
  \item Siapkan file konfigurasi contoh:
\begin{lstlisting}[language=bash, caption={Membuat file konfigurasi contoh}]
echo "PORT=8080" > ~/praktikum-os/week07-bash/sample-app.conf
cat ~/praktikum-os/week07-bash/sample-app.conf
\end{lstlisting}

  \item Tambahkan fungsi ke \texttt{.bashrc}:
\begin{lstlisting}[language=bash, caption={Menambahkan fungsi backup\_conf ke .bashrc}]
cat <<'EOF' >> ~/.bashrc

#: Praktikum Fungsi Shell ---
backup_conf() {
  if [ $# -ne 1 ]; then
    echo "Usage: backup_conf <file>"
    return 1
  fi

  local src="$1"
  local dst="$HOME/praktikum-os/week07-bash/backup"

  if [ ! -f "$src" ]; then
    echo "File tidak ditemukan: $src"
    return 2
  fi

  mkdir -p "$dst"
  cp -- "$src" "$dst/$(basename "$src").$(date +%F-%H%M%S).bak"
  echo "Backup selesai di $dst"
}
#: End Praktikum Fungsi Shell ---

EOF

source ~/.bashrc
\end{lstlisting}

  \item Uji fungsi:
\begin{lstlisting}[language=bash, caption={Menguji fungsi backup\_conf}]
backup_conf ~/praktikum-os/week07-bash/sample-app.conf
ls -lah ~/praktikum-os/week07-bash/backup
type backup_conf
\end{lstlisting}
\end{enumerate}

%% ------------------------------------------------------------------------
\section{Penyelesaian (Completion) dan History Bash}
\label{sec:completion-history}

Bash membantu pengguna mengetik lebih cepat melalui fitur \textit{completion}.
Saat Anda menekan tombol \texttt{Tab}, Bash mencoba melengkapi nama file,
direktori, atau perintah yang sedang diketik. Fitur ini terlihat sederhana,
tetapi dalam praktik administrasi sistem sangat membantu karena dapat
mengurangi kesalahan pengetikan path atau nama command.

Selain itu, Bash juga menyimpan \textit{history}, yaitu daftar perintah yang pernah
dijalankan. History berguna untuk mengulang langkah kerja, menelusuri perintah
yang pernah dipakai saat troubleshooting, dan mendokumentasikan workflow kerja.
Fitur pencarian history dengan \texttt{Ctrl+r} adalah salah satu senjata penting
administrator yang sering bekerja di terminal.

\subsection*{Fungsi Perintah yang Digunakan}

\begin{itemize}
  \item tombol \texttt{Tab}: completion nama file/perintah.
  \item \texttt{history}: melihat riwayat perintah.
  \item \texttt{history | grep pola}: mencari perintah tertentu dari history.
  \item \texttt{!\textless nomor\textgreater}: menjalankan ulang perintah berdasarkan nomor history.
  \item \texttt{!!}: menjalankan ulang perintah terakhir.
  \item \texttt{Ctrl+r}: pencarian interaktif ke history.
\end{itemize}

\begin{notebox}
History sangat membantu, tetapi juga berbahaya jika dipakai tanpa berpikir.
Menjalankan ulang perintah dengan \texttt{!!} atau \texttt{!\textless nomor\textgreater}
dapat berisiko jika perintah sebelumnya bersifat destruktif.
\end{notebox}

\subsection*{Praktikum 6.9: Menggunakan Completion Dasar dan Melihat History }

Tujuan: menggunakan fitur \texttt{Tab} dan history untuk mempercepat kerja.

\begin{enumerate}
  \item Pastikan file contoh tersedia:
\begin{lstlisting}[language=bash, caption={Menyiapkan file untuk completion}]
cd ~/praktikum-os/week07-bash/sampel
touch laporan-harian.log laporan-mingguan.log laporan-bulanan.log
ls
\end{lstlisting}

  \item Uji completion file:
  \begin{enumerate}
    \item Ketik \texttt{cat lap} lalu tekan \texttt{Tab} dua kali.
    \item Amati daftar file yang memiliki prefix \texttt{lap}.
    \item Ketik lebih spesifik, misalnya \texttt{cat laporan-h} lalu tekan \texttt{Tab}.
  \end{enumerate}

  \item Jalankan beberapa perintah sederhana:
\begin{lstlisting}[language=bash, caption={Menjalankan beberapa perintah untuk history}]
pwd
ls -lah
date
whoami
history | tail -n 10
\end{lstlisting}
\end{enumerate}

\subsection*{Praktikum 6.10: Menelusuri Perintah Diagnostik dengan History }

Tujuan: menggunakan history untuk mengulang langkah diagnosis sistem secara cepat.

\begin{enumerate}
  \item Jalankan beberapa perintah diagnostik:
\begin{lstlisting}[language=bash, caption={Menjalankan perintah diagnostik}]
df -h
free -h
uptime
ps aux | head
\end{lstlisting}

  \item Cari ulang perintah diagnostik dari history:
\begin{lstlisting}[language=bash, caption={Mencari perintah dari history}]
history | grep -E 'df -h|free -h|uptime|ps aux'
\end{lstlisting}

  \item Jalankan ulang salah satu perintah berdasarkan nomor history:
\begin{lstlisting}[language=bash, caption={Menjalankan ulang perintah dari history}]
!<NOMOR_HISTORY_ANDA>
\end{lstlisting}

  \item Simpan potongan history ke file dokumentasi:
\begin{lstlisting}[language=bash, caption={Menyimpan history ke file}]
history | tail -n 20 > ~/praktikum-os/week07-bash/diag-history.txt
cat ~/praktikum-os/week07-bash/diag-history.txt
\end{lstlisting}
\end{enumerate}

\begin{tipbox}
Jika paket \texttt{bash-completion} tersedia pada sistem Anda, maka completion
untuk banyak command akan menjadi lebih kaya. Namun, completion dasar untuk
nama file, direktori, dan command sudah tersedia tanpa konfigurasi tambahan.
\end{tipbox}

%% ------------------------------------------------------------------------
\section{Wildcards dan Ekspansi Nama File}
\label{sec:wildcards-ekspansi}

Bash mendukung \textit{wildcards} atau \textit{globbing} untuk mencocokkan banyak
nama file sekaligus. Fitur ini sangat penting untuk pekerjaan operasional karena
administrator sering harus memproses sekelompok file log, file backup, atau file
konfigurasi yang namanya berpola serupa.

Selain wildcard, Bash juga memiliki ekspansi lain seperti brace expansion dan tilde
expansion. Memahami perbedaan di antara semuanya membantu Anda bekerja lebih
cepat, tetapi juga lebih aman, terutama ketika suatu pola akan digunakan bersama
perintah seperti \texttt{cp}, \texttt{mv}, atau \texttt{rm}.

\subsection*{Fungsi Perintah yang Digunakan}

\begin{itemize}
  \item \texttt{*}: mencocokkan nol atau lebih karakter.
  \item \texttt{?}: mencocokkan tepat satu karakter.
  \item \texttt{[abc]} atau \texttt{[0-9]}: mencocokkan karakter tertentu/rentang.
  \item \texttt{\{a,b,c\}}: brace expansion.
  \item \texttt{\textasciitilde}: mengembang menjadi home directory pengguna.
  \item \texttt{echo pola}: mem-preview hasil ekspansi sebelum dipakai pada command berisiko.
\end{itemize}

\begin{notebox}
Kebiasaan aman saat bekerja dengan wildcard adalah mem-preview dulu hasilnya
menggunakan \texttt{echo}. Jangan langsung memakai \texttt{rm *.log} atau
\texttt{mv *.conf ...} tanpa melihat apa saja yang akan cocok dengan pola tersebut.
\end{notebox}

\subsection*{Praktikum 6.11: Mencoba Wildcard Dasar }

Tujuan: memahami pola wildcard dan ekspansi nama file.

\begin{enumerate}
  \item Masuk ke direktori sampel:
\begin{lstlisting}[language=bash, caption={Masuk ke direktori sampel}]
cd ~/praktikum-os/week07-bash/sampel
ls
\end{lstlisting}

  \item Coba beberapa pola wildcard:
\begin{lstlisting}[language=bash, caption={Mencoba wildcard dasar}]
ls *.log
ls catatan-?.txt
ls backup-0[12].tar
\end{lstlisting}

  \item Coba beberapa ekspansi lain:
\begin{lstlisting}[language=bash, caption={Mencoba brace dan tilde expansion}]
echo log-{pagi,siang,malam}.txt
echo ~
echo ~/praktikum-os/week04-bash
\end{lstlisting}
\end{enumerate}

\subsection*{Praktikum 6.12: Mengarsipkan Banyak Log Sekaligus }

Tujuan: menggunakan wildcard secara aman untuk pekerjaan rutin administrator.

\begin{enumerate}
  \item Siapkan file log tambahan:
\begin{lstlisting}[language=bash, caption={Menyiapkan file log tambahan}]
cd ~/praktikum-os/week07-bash/logs
touch access-01.log access-02.log access-03.log
ls
\end{lstlisting}

  \item Preview file yang akan diproses:
\begin{lstlisting}[language=bash, caption={Preview hasil wildcard}]
echo *.log
echo access-0?.log
\end{lstlisting}

  \item Pindahkan semua file log ke folder arsip:
\begin{lstlisting}[language=bash, caption={Memindahkan file log ke folder arsip}]
mkdir -p arsip-log
mv *.log arsip-log/
ls arsip-log
\end{lstlisting}

  \item Kompres folder arsip:
\begin{lstlisting}[language=bash, caption={Membuat arsip terkompresi}]
tar -czf arsip-log-$(date +%F).tar.gz arsip-log
ls -lah
\end{lstlisting}
\end{enumerate}

%% ------------------------------------------------------------------------
\section{Quoting dan Escaping di Bash}
\label{sec:quoting-escaping}

Salah satu sumber kesalahan yang paling sering terjadi di Bash adalah masalah
\textit{quoting}. Banyak pemula menganggap shell hanya membaca kata per kata,
padahal Bash melakukan banyak jenis ekspansi sebelum perintah dijalankan:
ekspansi variabel, wildcard, command substitution, dan pemisahan berdasarkan spasi.
Karena itu, penulisan tanda kutip yang tepat sangat penting.

Secara umum:
\begin{itemize}
  \item \textbf{Single quote} (\texttt{'...'}) menjaga isi tetap literal.
  \item \textbf{Double quote} (\texttt{"..."}) tetap mengizinkan ekspansi variabel
        dan command substitution, tetapi mempertahankan spasi sebagai satu argumen.
  \item \textbf{Backslash} (\texttt{\textbackslash}) digunakan untuk meng-escape
        satu karakter tertentu.
\end{itemize}

Dalam pekerjaan administrator, quoting yang benar sangat penting saat menangani
nama file yang mengandung spasi, tanda kurung siku, tanda bintang, atau saat
membangun command berbasis variabel.

\subsection*{Fungsi Perintah yang Digunakan}

\begin{itemize}
  \item \texttt{echo '...'}: menampilkan teks literal.
  \item \texttt{echo "..."}: menampilkan teks dengan ekspansi variabel.
  \item \texttt{\textbackslash}: meng-escape karakter khusus.
  \item \texttt{cp -- "sumber" "tujuan"}: menyalin file dengan nama yang kompleks.
  \item \texttt{printf '\%s\textbackslash n' "\$var"}: menampilkan nilai variabel dengan aman.
\end{itemize}

\begin{tipbox}
Aturan praktis yang sangat berguna: \textbf{selalu beri tanda kutip ganda pada variabel
yang merepresentasikan path atau nama file}, misalnya \texttt{"\$file"}.
\end{tipbox}

\subsection*{Praktikum 6.13: Membedakan Single Quote, Double Quote, dan Escape }

Tujuan: memahami efek tanda kutip terhadap ekspansi Bash.

\begin{enumerate}
  \item Uji single quote dan double quote:
\begin{lstlisting}[language=bash, caption={Menguji perbedaan jenis quote}]
echo '$USER bekerja di $HOME'
echo "$USER bekerja di $HOME"
\end{lstlisting}

  \item Uji escape karakter spasi:
\begin{lstlisting}[language=bash, caption={Menggunakan escape pada path}]
cd ~/praktikum-os/week07-bash/ruang-nama
ls laporan\ server\ april.txt
\end{lstlisting}

  \item Uji akses file yang sama dengan double quote:
\begin{lstlisting}[language=bash, caption={Menggunakan double quote pada path}]
cat "laporan server april.txt"
\end{lstlisting}
\end{enumerate}

\subsection*{Praktikum 6.14: Menangani File dengan Nama Sulit Secara Aman }

Tujuan: memproses file yang mengandung spasi dan karakter khusus tanpa error.

\begin{enumerate}
  \item Pastikan file target tersedia:
\begin{lstlisting}[language=bash, caption={Memastikan file target tersedia}]
cd ~/praktikum-os/week07-bash/ruang-nama
ls -lah
\end{lstlisting}

  \item Salin file dengan nama kompleks ke folder backup:
\begin{lstlisting}[language=bash, caption={Menyalin file dengan nama kompleks secara aman}]
cp -- "backup [mingguan] server.conf" \
  "$HOME/praktikum-os/week07-bash/backup/backup-mingguan-server.conf"
\end{lstlisting}

  \item Gunakan variabel untuk memproses path dengan aman:
\begin{lstlisting}[language=bash, caption={Menggunakan variabel berisi path dengan aman}]
file_asli="$HOME/praktikum-os/week07-bash/ruang-nama/backup [mingguan] server.conf"
file_salinan="$HOME/praktikum-os/week07-bash/backup/backup-mingguan-server-v2.conf"

cp -- "$file_asli" "$file_salinan"
ls -lah "$HOME/praktikum-os/week07-bash/backup"
\end{lstlisting}

  \item Tampilkan daftar file hasil backup:
\begin{lstlisting}[language=bash, caption={Menampilkan file hasil backup dengan aman}]
for file in "$HOME"/praktikum-os/week07-bash/backup/*; do
  printf 'Hasil backup: %s\n' "$file"
done
\end{lstlisting}
\end{enumerate}

\begin{notebox}
Kesalahan klasik adalah menulis \texttt{cp \$file tujuan/} tanpa tanda kutip.
Jika isi variabel \texttt{\$file} mengandung spasi, Bash akan memecahnya menjadi
beberapa argumen dan command bisa gagal atau bahkan memproses file yang salah.
\end{notebox}

%% ------------------------------------------------------------------------
\section{Tugas Praktikum}
\label{sec:tugas-praktikum-bash}

\begin{warningbox}
Seluruh tugas dikerjakan pada workspace praktikum sendiri, misalnya
\texttt{\$HOME/praktikum-os/week04-bash}, kecuali jika instruksi secara eksplisit
meminta membaca file sistem seperti \texttt{/etc} atau \texttt{/proc}. Hindari
mengubah file sistem penting secara langsung tanpa backup.
\end{warningbox}

\subsection*{Tugas Praktikum 1: Toolkit Bash Administrator Pribadi}
\textbf{Konteks riil:} seorang administrator sering mengulang perintah yang sama
setiap hari. Agar pekerjaan lebih efisien dan konsisten, ia perlu memiliki
\textit{toolkit} Bash pribadi yang otomatis aktif setiap login.

\textbf{Instruksi tugas:}
\begin{enumerate}
    \item Tambahkan konfigurasi pada \texttt{.bashrc} untuk:
    \begin{itemize}
        \item menambahkan direktori \texttt{bin} pribadi ke \texttt{PATH},
        \item membuat minimal 2 alias yang membantu kerja harian,
        \item membuat minimal 1 fungsi shell yang berguna untuk administrasi.
    \end{itemize}
    \item Pastikan konfigurasi tersebut aktif kembali saat membuka shell login.
    \item Buat satu script sederhana di direktori \texttt{bin} pribadi, misalnya
    script untuk menampilkan ringkasan sistem.
    \item Uji dari direktori yang berbeda untuk memastikan script dapat dipanggil
    tanpa menuliskan path lengkap.
    \item Simpan bukti pengujian ke file \texttt{toolkit-bash-report.txt}.
\end{enumerate}

\textbf{Minimal luaran:}
\begin{itemize}
    \item isi blok konfigurasi yang ditambahkan ke \texttt{.bashrc},
    \item output \texttt{echo \$PATH},
    \item output \texttt{type} untuk alias, fungsi, dan script,
    \item file laporan \texttt{toolkit-bash-report.txt}.
\end{itemize}

\subsection*{Tugas Praktikum 2: Audit File Konfigurasi dan Logging Aman}
\textbf{Konteks riil:} saat troubleshooting, administrator sering perlu
menginventarisasi file konfigurasi dan memisahkan output normal dari pesan error.

\textbf{Instruksi tugas:}
\begin{enumerate}
    \item Buat file laporan bernama
    \texttt{audit-konfigurasi-\$(date +\%F).txt}.
    \item Cari file \texttt{*.conf} di dalam \texttt{/etc} dan simpan hasilnya
    ke file laporan.
    \item Catat jumlah total file konfigurasi yang ditemukan.
    \item Jika ada pesan error, simpan ke file terpisah, misalnya
    \texttt{audit-error.log}.
    \item Tampilkan isi laporan ke terminal dan sekaligus simpan menggunakan
    \texttt{tee}.
    \item Tambahkan ringkasan singkat 3--5 baris yang menjelaskan mengapa
    pemisahan \texttt{stdout} dan \texttt{stderr} penting dalam audit sistem.
\end{enumerate}

\textbf{Syarat konsep yang harus muncul:}
\begin{itemize}
    \item redirection \texttt{>}, \texttt{2>}, atau \texttt{\&>},
    \item pipeline,
    \item \texttt{tee},
    \item penggunaan variabel atau \texttt{command substitution}.
\end{itemize}

\textbf{Minimal luaran:}
\begin{itemize}
    \item file laporan audit,
    \item file log error,
    \item perintah yang digunakan,
    \item analisis singkat hasil audit.
\end{itemize}

\subsection*{Tugas Praktikum 3: Mini Health Check Harian Server}
\textbf{Konteks riil:} administrator perlu membuat pemeriksaan cepat
(\textit{health check}) untuk mengetahui kondisi dasar server sebelum dan sesudah
maintenance.

\textbf{Instruksi tugas:}
\begin{enumerate}
    \item Buat script Bash bernama \texttt{daily-healthcheck} pada direktori
    \texttt{bin} pribadi.
    \item Script minimal harus menampilkan:
    \begin{itemize}
        \item tanggal dan waktu,
        \item hostname,
        \item user aktif,
        \item shell aktif,
        \item uptime,
        \item penggunaan memori,
        \item penggunaan filesystem root,
        \item 10 baris terakhir history command yang relevan dengan pengecekan.
    \end{itemize}
    \item Simpan hasil ke file log harian, misalnya
    \texttt{healthcheck-\$(date +\%F).log}.
    \item Tampilkan hasil ke terminal dan ke file secara bersamaan.
    \item Jika Anda menggunakan pipeline dengan \texttt{tee}, cek juga status
    exit command utama.
\end{enumerate}

\textbf{Syarat konsep yang harus muncul:}
\begin{itemize}
    \item environment variable,
    \item \texttt{PATH},
    \item alias atau fungsi pendukung,
    \item history,
    \item \texttt{tee},
    \item penanganan error dasar.
\end{itemize}

\textbf{Minimal luaran:}
\begin{itemize}
    \item file script yang executable,
    \item contoh isi file log hasil eksekusi,
    \item penjelasan singkat fungsi tiap bagian script.
\end{itemize}

\subsection*{Tugas Praktikum 4: Penanganan File dengan Nama Kompleks dan Arsip Aman}
\textbf{Konteks riil:} file hasil backup, ekspor, atau laporan sering memiliki nama
yang mengandung spasi atau karakter khusus. Administrator harus tetap dapat
memproses file-file tersebut tanpa salah target.

\textbf{Instruksi tugas:}
\begin{enumerate}
    \item Buat minimal 4 file contoh dengan nama yang bervariasi, termasuk:
    \begin{itemize}
        \item nama file yang mengandung spasi,
        \item nama file yang mengandung tanda kurung siku atau karakter khusus,
        \item file dengan pola nama serupa untuk diuji dengan wildcard.
    \end{itemize}
    \item Tunjukkan perbedaan hasil jika file diakses tanpa quoting dan dengan
    quoting yang benar.
    \item Lakukan \textit{preview} wildcard dengan \texttt{echo} sebelum dipakai
    untuk operasi nyata.
    \item Salin file-file tersebut ke direktori backup dengan nama yang aman.
    \item Buat arsip \texttt{tar.gz} dari hasil backup.
    \item Simpan riwayat perintah yang Anda gunakan ke file
    \texttt{riwayat-arsip.txt}.
\end{enumerate}

\textbf{Syarat konsep yang harus muncul:}
\begin{itemize}
    \item single quote, double quote, dan escaping,
    \item wildcard,
    \item variabel path,
    \item history,
    \item operasi file lanjutan yang aman.
\end{itemize}

\textbf{Minimal luaran:}
\begin{itemize}
    \item daftar file awal,
    \item daftar file hasil backup,
    \item file arsip \texttt{tar.gz},
    \item file \texttt{riwayat-arsip.txt},
    \item refleksi singkat tentang pentingnya quoting di Bash.
\end{itemize}

%% ------------------------------------------------------------------------
\section{Rangkuman}
\label{sec:rangkuman-bash}

\begin{itemize}
  \item Bash adalah shell yang sangat umum digunakan pada Linux dan menjadi alat utama
        interaksi administrator dengan sistem.
  \item \texttt{.bashrc} dipakai untuk konfigurasi shell interaktif, sedangkan
        \texttt{.bash\_profile} dipakai pada konteks login shell.
  \item Variabel lingkungan dan \texttt{PATH} menentukan konteks kerja shell serta
        bagaimana Bash menemukan command dan script.
  \item Alias cocok untuk shortcut sederhana, sedangkan fungsi shell cocok untuk
        logika yang lebih kompleks dan berulang.
  \item Completion dan history membantu mempercepat kerja serta mengurangi
        kesalahan pengetikan.
  \item Wildcard dan ekspansi nama file memudahkan pemrosesan banyak file sekaligus,
        tetapi harus digunakan dengan kebiasaan preview yang aman.
  \item Quoting dan escaping adalah fondasi penting agar Bash menangani variabel,
        spasi, dan karakter khusus dengan benar.
\end{itemize}

%% ------------------------------------------------------------------------
\section{Referensi dan Bahan Belajar Lanjutan}
\label{sec:referensi-lanjutan-bash}

Untuk pendalaman materi, pembaca disarankan mempelajari referensi berikut:

\begin{enumerate}
    \item GNU Bash Reference Manual.
    \item Bash manual page (\texttt{man bash}).
    \item Bash Guide for Beginners.
    \item William Shotts, \textit{The Linux Command Line}.
    \item Evi Nemeth, Garth Snyder, Trent R. Hein, Ben Whaley, Dan Mackin,
    \textit{UNIX and Linux System Administration Handbook}, 5th Edition.
    \item Arnold Robbins dan Nelson H. F. Beebe,
    \textit{Classic Shell Scripting}.
    \item Mark G. Sobell,
    \textit{A Practical Guide to Linux Commands, Editors, and Shell Programming}.
    \item Richard Blum dan Christine Bresnahan,
    \textit{Linux Command Line and Shell Scripting Bible}.
\end{enumerate}